\documentclass[12pt, a4paper]{article}
\usepackage{amsmath, amsthm, amssymb}
\usepackage{geometry}
\usepackage{microtype}
\usepackage{enumitem}
\geometry{margin=1in}

% ── Theorem environments ──────────────────────────────────────────────────────
\theoremstyle{plain}
\newtheorem{lemma}{Lemma}
\newtheorem{theorem}{Theorem}

\theoremstyle{definition}
\newtheorem{axiomentry}{Axiom}
\newtheorem{defentry}{Definition}

\title{\textbf{Peano Arithmetic}}
\author{0xHaru}
\date{\today}

\begin{document}
\maketitle

% ─────────────────────────────────────────────────────────────────────────────
\section*{Peano Axioms}
% ─────────────────────────────────────────────────────────────────────────────

We first define the \textbf{successor function} $S : \mathbb{N} \to \mathbb{N}$
by $S(n) = n + 1$. For example, $S(0) = 1$.

The four Peano axioms are:

\begin{axiomentry}[A1]
$0$ is a natural number.
\end{axiomentry}

\begin{axiomentry}[A2]
For every natural number $n$, $S(n)$ is a natural number.
\end{axiomentry}

\begin{axiomentry}[A3]
For all natural numbers $n$ and $m$, if $S(n) = S(m)$ then $n = m$.
\end{axiomentry}

\begin{axiomentry}[A4]
There does not exist a natural number $n$ such that $S(n) = 0$.
\end{axiomentry}

\subsection*{Addition}

Addition on $\mathbb{N}$ is defined by the following two rules,
for all $n, m \in \mathbb{N}$:

\begin{defentry}[D1]
$n + 0 = n$
\end{defentry}

\begin{defentry}[D2]
$n + S(m) = S(n + m)$
\end{defentry}

% ─────────────────────────────────────────────────────────────────────────────
\section*{Lemma 1: $S(a) = a + 1$}
% ─────────────────────────────────────────────────────────────────────────────

\begin{lemma}\label{lem:succ}
For all $a \in \mathbb{N}$, $\;S(a) = a + 1$.
\end{lemma}

\begin{proof}
\[
  S(a)
    \;\underset{\mathrm{D1}}{=}\; S(a + 0)
    \;\underset{\mathrm{D2}}{=}\; a + S(0)
    \;\underset{\mathrm{Def.\,of\,1}}{=}\; a + 1.
\]
\end{proof}

% ─────────────────────────────────────────────────────────────────────────────
\section*{Proof that $2 \neq 1$}
% ─────────────────────────────────────────────────────────────────────────────

\begin{theorem}
$2 \neq 1$, i.e., $S(S(0)) \neq S(0)$.
\end{theorem}

\begin{proof}
We want to prove $\neg\bigl(S(S(0)) = S(0)\bigr)$.
Suppose for contradiction that $S(S(0)) = S(0)$ is true. Then:
\[
  S(S(0)) = S(0) \;\underset{\mathrm{A3}}{\implies}\; S(0) = 0.
\]
But A4 states that there is no natural number $n$ with $S(n) = 0$,
so $S(0) = 0$ is impossible. Contradiction.
\end{proof}

% ─────────────────────────────────────────────────────────────────────────────
\section*{Lemma 2: $S(a) + b = S(a + b)$}
% ─────────────────────────────────────────────────────────────────────────────

\begin{lemma}\label{lem:succ-add}
For all $a, b \in \mathbb{N}$, $\;S(a) + b = S(a + b)$.
\end{lemma}

\begin{proof}
By induction on $b$.

\medskip
\noindent\textbf{Base case} ($b = 0$):
\[
\begin{aligned}
  S(a) + 0 &\;=\; S(a + 0)\\
  S(a) &\;=\; S(a) &&\text{by D1 (twice)}\\
  a &\;=\; a &&\text{by A3.}
\end{aligned}
\]

\medskip
\noindent\textbf{Induction Hypothesis (I.H.):} $S(a) + b = S(a + b)$.

\medskip
\noindent\textbf{Induction Step.} We show $S(a) + S(b) = S(a + S(b))$:
\[
\begin{aligned}
  S(a) + S(b) &\;=\; S\bigl(a + S(b)\bigr)\\
  S\bigl(S(a) + b\bigr) &\;=\; S\bigl(a + S(b)\bigr) &&\text{by D2}\\
  S\bigl(S(a) + b\bigr) &\;=\; S\bigl(S(a + b)\bigr) &&\text{by D2}\\
  S\bigl(S(a + b)\bigr) &\;=\; S\bigl(S(a + b)\bigr) &&\text{by I.H.}\\
  a + b &\;=\; a + b &&\text{by A3 (twice).}
\end{aligned}
\]
\qedhere
\end{proof}

% ─────────────────────────────────────────────────────────────────────────────
\section*{Proof of Associativity}
% ─────────────────────────────────────────────────────────────────────────────

\begin{theorem}
For all $a, b, c \in \mathbb{N}$,
\[
  (a + b) + c \;=\; a + (b + c).
\]
\end{theorem}

\begin{proof}
By induction on $c$.

\medskip
\noindent\textbf{Base case} ($c = 0$):
\[
  (a + b) + 0
  \;\underset{\mathrm{D1}}{=}\; a + b
  \;\underset{\mathrm{D1}}{=}\; a + (b + 0).
\]

\medskip
\noindent\textbf{Induction Hypothesis (I.H.):}
$(a + b) + c = a + (b + c)$.

\medskip
\noindent\textbf{Induction Step.} We show
$(a + b) + S(c) = a + (b + S(c))$:
\[
\begin{aligned}
  (a + b) + S(c)
    &\;=\; S\bigl((a + b) + c\bigr) &&\text{by D2}\\
    &\;=\; S\bigl(a + (b + c)\bigr)   &&\text{by I.H.}\\
    &\;=\; a + S(b + c)               &&\text{by D2}\\
    &\;=\; a + (b + S(c))             &&\text{by D2.}
\end{aligned}
\]
The I.H.\ holds for $S(c)$. \qedhere
\end{proof}

% ─────────────────────────────────────────────────────────────────────────────
\section*{Proof of the Identity Element}
% ─────────────────────────────────────────────────────────────────────────────

\begin{theorem}
For all $a \in \mathbb{N}$, $\;0 + a = a$.
\end{theorem}

\begin{proof}
D1 states that $0$ is a \emph{right} identity ($a + 0 = a$).
We prove that $0$ is also a \emph{left} identity by induction on $a$.

\medskip
\noindent\textbf{Base case} ($a = 0$):
$\;0 + 0 = 0\;$ by D1.

\medskip
\noindent\textbf{Induction Hypothesis (I.H.):} $0 + a = a$.

\medskip
\noindent\textbf{Induction Step.} We show $0 + S(a) = S(a)$:
\[
\begin{aligned}
  0 + S(a)
    &\;=\; S(0 + a) &&\text{by D2}\\
    &\;=\; S(a)     &&\text{by I.H.}
\end{aligned}
\]
\qedhere
\end{proof}

% ─────────────────────────────────────────────────────────────────────────────
\section*{Proof of Commutativity}
% ─────────────────────────────────────────────────────────────────────────────

\begin{theorem}
For all $a, b \in \mathbb{N}$, $\;a + b = b + a$.
\end{theorem}

\begin{proof}
By induction on $b$. We first establish the two base cases $b = 0$ and
$b = S(0) = 1$ (i.e.\ we prove that $0$ and $1$ commute with everything).

% ---------- base case b = 0 ----------
\medskip
\noindent\textbf{Base case} ($b = 0$):
\[
  a + 0
  \;\underset{\mathrm{D1}}{=}\; a
  \;\underset{\text{Identity}}{=}\; 0 + a.
\]
The second equality is the left-identity theorem proved above.

% ---------- base case b = 1 ----------
\medskip
\noindent\textbf{Base case} ($b = 1$):
We prove $a + 1 = 1 + a$ by induction on $a$
\emph{(an induction proof within an induction proof)}.

\begin{itemize}[leftmargin=2.5em]
  \item \textbf{Inner base case} ($a = 0$):
  \[
  \begin{aligned}
    0 + 1 &\;=\; 1 + 0 \\
    1     &\;=\; 1 + 0 &&\text{by the Identity Element}\\
    1     &\;=\; 1     &&\text{by D1.}
  \end{aligned}
  \]

  \item \textbf{Inner Induction Hypothesis (I.H.):} $a + 1 = 1 + a$.

  \item \textbf{Inner Induction Step.}
  We show $S(a) + 1 = 1 + S(a)$:
  \[
  \begin{aligned}
    S(a) + 1
      &\;=\; S(a) + S(0) &&\text{by definition of }1\\
      &\;=\; S(S(a) + 0) &&\text{by D2}\\
      &\;=\; S(S(a)    ) &&\text{by D1}\\
      &\;=\; S(a + 1)    &&\text{by Lemma~\ref{lem:succ}}\\
      &\;=\; S(1 + a)    &&\text{by I.H.}\\
      &\;=\; 1 + S(a)    &&\text{by D2.}
  \end{aligned}
  \]
\end{itemize}

We have proved the base case $b = 1$.

% ---------- induction hypothesis & step ----------
\medskip
\noindent\textbf{Induction Hypothesis (I.H.):} $a + b = b + a$.

\medskip
\noindent\textbf{Induction Step.} We show $a + S(b) = S(b) + a$:
\[
\begin{aligned}
  a + S(b)
    &\;=\; a + (b + 1)   &&\text{by Lemma~\ref{lem:succ}}\\
    &\;=\; (a + b) + 1   &&\text{by Associativity}\\
    &\;=\; (b + a) + 1   &&\text{by I.H.}\\
    &\;=\; b + (a + 1)   &&\text{by Associativity}\\
    &\;=\; b + (1 + a)   &&\text{by base case }b = 1\\
    &\;=\; (b + 1) + a   &&\text{by Associativity}\\
    &\;=\; S(b) + a      &&\text{by Lemma~\ref{lem:succ}.}
\end{aligned}
\]
\qedhere
\end{proof}

\end{document}
